% Created 2022-03-08 Tue 12:01
% Intended LaTeX compiler: pdflatex
\documentclass[11pt]{article}
\usepackage[utf8]{inputenc}
\usepackage[T1]{fontenc}
\usepackage{graphicx}
\usepackage{longtable}
\usepackage{wrapfig}
\usepackage{rotating}
\usepackage[normalem]{ulem}
\usepackage{amsmath}
\usepackage{amssymb}
\usepackage{capt-of}
\usepackage{hyperref}
\usepackage{natbib}
\usepackage{siunitx}

\author{Chiara P. Montagna, INGV Pisa}
\title{Slug ascent at Stromboli.}
\hypersetup{
 pdfauthor={Chiara P. Montagna, INGV Pisa},
 pdftitle={Modeling slug rise at Stromboli, and the associated VLP seismic signal},
 pdfkeywords={},
 pdfsubject={},
 pdfcreator={Emacs 26.3 (Org mode 9.5.2)}, 
 pdflang={English}}
\begin{document}

\bibliographystyle{plainnat} 
\bibpunct{(}{)}{;}{a}{,}{,}

\maketitle

\section{Method}
The ascent of a gas slug at Stromboli is modeled using
\href{https://zenodo.org/record/5031825\#.YgP13vso89k}{MagmaFOAM}.
Ground deformation signals generated by slug ascent are calculated using
\href{https://zenodo.org/record/4557510\#.YgP3svso89k}{syntheticSignals},
and compared to seismic and strain records in the VLP band \citep{Chouet2003,Mattia2021}.

Pressure at the boundaries of the fluid system is used to calculate
ground displacement at the Earth's surface \citep{Longo2012bv}.

\section{Slug ascent}
The conduit is \SI{170}{m} long and \SI{30}{m} wide. Gas
  enters from the bottom, forming Taylor bubbles that rise towards the top
  (\href{stromboli/movies/slug3.ogv}{movie}).

\section{Ground deformation}


\bibliography{/home/montagna/Dropbox/biblio.bib}

\end{document}